\section{Introduction}
\label{sec:introduction}

Flash floods are the dominant natural hazard in the Western Himalaya,
accounting for the majority of disaster-related fatalities and economic losses
across Himachal Pradesh (HP), Uttarakhand, and Jammu and Kashmir each monsoon season
\citep{kumar2022hp, dixit2026nhess}.
The 2023 HP monsoon season was particularly catastrophic: 404 confirmed fatalities,
more than 4,700 roads damaged or blocked, and direct economic losses exceeding
Rs.\ 9,905 crore (\$1.2 billion USD), making it one of the three most destructive
flood seasons in Indian recorded history \citep{ndma2023, hiflodot2025}.
These losses were concentrated along the Beas, Satluj, and Chenab valleys,
where steep terrain, high seasonal precipitation, and dense infrastructure create
exceptional compound hazard conditions.

Flash flood susceptibility mapping (FFSM) provides the spatial risk baseline
required for land-use planning, infrastructure protection, early warning system design,
and emergency response allocation.
Machine-learning (ML) methods have displaced traditional frequency-ratio and bivariate
statistical approaches since approximately 2018
\citep{abedi2021, bentivoglio2022review, ghosh2023karnali},
with ensemble tree models (Random Forest, XGBoost, LightGBM) and their
stacked combinations now representing the state of the art in AUC performance
\citep{chen2023rules, kosi2025ml}.
Despite this methodological progress, critical gaps persist,
particularly in the Western Himalayan context.

\subsection{Research Gaps}
\label{sec:gaps}

\paragraph{Gap 1 --- No comprehensive ML study for HP.}
Only one ML paper targets HP specifically:
\citet{saha2023} mapped susceptibility for the Beas basin using
Multivariate Adaptive Regression Splines (MARS) and RF, achieving AUC\,=\,0.88
in a book chapter published with Springer.
No journal article covers the full state or the climatologically distinct upper basins
(Lahaul-Spiti, Kinnaur, upper Chamba).
A CNN U-Net study (Research Square, 2025) remains a preprint.
The absence of a comprehensive benchmark covering all HP basins represents a clear
publication gap.

\paragraph{Gap 2 --- Spatial autocorrelation in validation.}
The overwhelming majority of susceptibility studies use random train-test splits,
which ignore spatial autocorrelation and inflate AUC by 5--15\%
\citep{valavi2019, roberts2017cross}.
Spatial block cross-validation -- where geographically contiguous units form the test fold --
has been advocated since at least 2017 but remains rarely applied in FFSM literature
for the Indian Himalayan region.

\paragraph{Gap 3 --- No uncertainty quantification.}
All published HP susceptibility maps produce a single probability estimate per pixel,
conveying no information about prediction confidence.
Decision-makers at HP SDMA are therefore unable to distinguish zones where
the model is highly certain from zones where uncertainty is large.
Conformal prediction \citep{angelopoulos2023conformal} provides
statistically guaranteed coverage intervals and is directly applicable to
pre-trained ML classifiers, yet no FFSM paper has applied it to produce
coverage-guaranteed susceptibility intervals.

\paragraph{Gap 4 --- Absence of graph-based spatial dependencies.}
Standard ML approaches treat every pixel or catchment as conditionally independent
given the conditioning factors, ignoring the hydrological reality that an
extreme-runoff event in an upstream catchment physically elevates flood risk
in all downstream catchments through increased channel discharge
\citep{bentivoglio2022review}.
Graph Neural Networks (GNNs) are ideally suited to encode this structure:
nodes represent sub-watersheds, directed edges encode river connectivity,
and message-passing aggregates upstream state into downstream risk predictions.
No published flash flood susceptibility paper applies GNNs as of early 2026.

\paragraph{Gap 5 --- Cloud-contaminated or absent flood inventories.}
Most HP studies rely on documentary records (NDMA, local government) that are
spatially imprecise and temporally incomplete.
Sentinel-1 Synthetic Aperture Radar (SAR) operates independently of cloud cover,
making it uniquely suitable for monsoon-season flood mapping
\citep{nagamani2024sar}.
A multi-year, SAR-derived, standardised flood inventory for HP has not been published.

\subsection{Objectives and Contributions}
\label{sec:objectives}

This study addresses all five gaps through the following specific contributions:

\begin{enumerate}
  \item \textbf{Multi-year SAR flood inventory.}
    We construct the first systematic, multi-temporal Sentinel-1 SAR flood inventory
    for HP (2018--2024), processed in Google Earth Engine using standardised
    change-detection methodology and combined with the HiFlo-DAT historical database.

  \item \textbf{Graph Neural Network susceptibility model.}
    We build a directed watershed connectivity graph for HP and train a
    GraphSAGE \citep{hamilton2017graphsage} model on it,
    capturing upstream--downstream flood propagation structure that
    pixel-based models cannot represent.
    This is, to our knowledge, the first application of GNNs to FFSM.

  \item \textbf{Conformal prediction intervals.}
    We apply split-conformal prediction (MAPIE; \citealt{taquet2022mapie})
    to the best-performing model, producing susceptibility maps with
    statistically guaranteed 90\% coverage intervals for the first time in FFSM.

  \item \textbf{Spatial block cross-validation.}
    All models are evaluated using leave-one-basin-out spatial block CV
    (five folds: Beas, Satluj, Chenab, Ravi, Yamuna),
    providing bias-corrected AUC, F1, and Cohen's Kappa estimates.

  \item \textbf{State-wide risk assessment with infrastructure overlay.}
    We produce the first comprehensive flash flood susceptibility map
    covering all five major HP river basins
    and quantify exposure of roads, bridges, hydroelectric installations,
    and tourist facilities to high-susceptibility zones.
\end{enumerate}

The remainder of this paper is organised as follows.
Section~\ref{sec:study_area} describes the study area.
Section~\ref{sec:data_methods} details data sources, preprocessing,
and the ML and GNN methods.
Section~\ref{sec:results} presents model performance, susceptibility maps,
conformal coverage analysis, and SHAP factor importance.
Section~\ref{sec:discussion} discusses the contribution of graph structure to accuracy,
uncertainty communication for SDMA, and study limitations.
Section~\ref{sec:conclusion} summarises key findings and policy implications.
